\section{Buenas prácticas de backend}

Estas pautas son una introducción independiente del lenguaje o la arquitectura que finalmente utilice AUTOVIDA. Se concretarán cuando el equipo acuerde la tecnología, los perfiles y los datos que realmente necesita el proyecto.

\subsection{Datos, validación y consistencia}
\begin{itemize}[leftmargin=*]
  \item Validar cada entrada en el servidor: presencia de campos requeridos, tipo, longitud y valores permitidos. Rechazar lo inválido con un mensaje controlado y comprensible para la interfaz.
  \item Aplicar las reglas de negocio en un lugar consistente. No duplicar reglas contradictorias entre pantallas y servidor.
  \item Guardar únicamente los datos necesarios para las funciones aprobadas. Usar nombres y formatos consistentes y establecer relaciones claras entre usuarios, tareas y actividades.
  \item Evitar estados imposibles o duplicados mediante restricciones y operaciones de escritura seguras. Si una operación debe completarse como un todo, impedir que quede guardada a medias.
  \item Definir cómo se corregirán errores de datos y cómo se harán copias de seguridad y restauraciones, según el alcance y la infraestructura aprobados.
\end{itemize}

\subsection{Permisos, seguridad y privacidad}
\begin{itemize}[leftmargin=*]
  \item Comprobar permisos en cada operación protegida, no solo en la interfaz. Aplicar el mínimo acceso necesario para cada perfil.
  \item No incluir contraseñas, tokens, claves ni datos personales innecesarios en el código, mensajes de error o registros técnicos. Las credenciales se gestionan mediante la configuración segura del entorno elegido.
  \item No devolver información de otras personas al consultar o modificar una tarea. Probar expresamente los casos de acceso permitido y denegado.
  \item Tratar los datos del alumnado con especial cuidado: limitar su uso al propósito acordado, restringir quién puede verlos y no reutilizarlos en pruebas sin autorización y protección adecuada.
  \item No mostrar trazas internas, consultas, rutas del servidor o detalles técnicos a las personas usuarias. Registrar solo lo necesario para investigar fallos y con acceso limitado.
\end{itemize}

\subsection{Errores y disponibilidad}
El servidor debe distinguir entre una operación completada, una entrada incorrecta, falta de permisos y un fallo inesperado. La interfaz recibirá una respuesta previsible y podrá explicarle a la persona qué ha pasado y qué puede hacer. Los fallos técnicos no deben borrar silenciosamente datos ni presentarse como si la tarea se hubiera guardado.

El equipo identificará las operaciones importantes, probará sus casos normales y de error y documentará las dependencias de las que dependen. No se prometerá disponibilidad, recuperación o seguridad absoluta: se acordarán medidas concretas, se probarán y se anotarán sus límites.

\subsection{Práctica guiada}
En grupos, recorrer el flujo ``crear una tarea''. Enumerar datos de entrada, reglas de validación, quién puede crearla, qué se guarda, respuestas posibles y qué mensaje muestra la interfaz si no se puede guardar. Revisar que no se confunda un fallo con una operación completada.
